\chapter{Mouth diseases and conditions}

\section*{Definition and Overview}
Mouth diseases and conditions, or oral pathologies, encompass a broad range of disorders affecting the oral mucosa, salivary glands, jaws, and tongue. These include infectious diseases (viral, bacterial, and fungal), autoimmune and inflammatory disorders (such as oral lichen planus and recurrent aphthous stomatitis), benign and malignant neoplasms (primarily oral squamous cell carcinoma), and salivary gland disorders (like sialadenitis or dry mouth). Because the oral cavity is highly vascular and sensitive, it frequently reflects underlying systemic illnesses (such as haematological malignancies, nutritional deficiencies, and autoimmune diseases), making the identification and management of oral pathologies critical for overall systemic health.

\section*{Epidemiology}
The epidemiology of mouth diseases is diverse, reflecting the varied nature of these conditions. Recurrent aphthous stomatitis (mouth ulcers) is the most common mucosal condition, affecting up to 20\% of the general population worldwide. Oral candidiasis is highly prevalent in specific cohorts, including newborns, denture wearers, and up to 90\% of patients with advanced HIV/AIDS. Oral lichen planus affects approximately 1--2\% of adults, predominantly females over the age of 40. Oral cancer, specifically oral squamous cell carcinoma, is a major global health concern, representing about 2--3\% of all cancers diagnosed in Australia. It has a strong male predilection and occurs primarily in individuals over the age of 50 with a history of heavy tobacco and alcohol use.

\section*{Aetiology and Pathophysiology}
Oral pathologies develop through diverse infectious, immunological, neoplastic, and environmental pathways.
\begin{itemize}
    \item \textbf{Infectious agents:} Pathogens including Herpes Simplex Virus (HSV-1, causing cold sores), Candida albicans (causing thrush), and Coxsackievirus (causing hand, foot, and mouth disease) directly invade and damage oral epithelial cells.
    \item \textbf{T-cell mediated autoimmunity:} In conditions like oral lichen planus, cytotoxic T-cells target and destroy basal keratinocytes in the oral epithelium, leading to hyperkeratotic or erosive mucosal lesions.
    \item \textbf{Carcinogen exposure:} Chronic exposure of the oral mucosa to tobacco carcinogens and ethanol (alcohol) induces genetic mutations in squamous cells, leading to dysplasia and progression to squamous cell carcinoma.
    \item \textbf{Systemic disease manifestations:} Autoimmune conditions (such as lupus or Crohn's disease) and nutritional deficiencies (such as iron or vitamin B-complex depletion) impair mucosal cell turnover, resulting in ulceration and glossitis.
    \item \textbf{Salivary gland obstruction and infection:} Calcification within salivary ducts (sialolithiasis) leads to salivary stasis, which predisposes the gland to secondary bacterial infection (sialadenitis), typically by Staphylococcus aureus.
\end{itemize}

\section*{Clinical Features}
Mouth diseases present with a wide spectrum of mucosal lesions, pain, and functional difficulties.
\begin{itemize}
    \item \textbf{Mucosal ulcerations:} Painful, well-demarcated erosions with an erythematous halo (aphthous ulcers), or chronic, non-healing ulcers with raised, indurated borders (concerning for malignancy).
    \item \textbf{White and red patches:} Non-wipeable, hyperkeratotic white plaques (leukoplakia) or erythematous patches (erythroplakia) on the tongue, floor of the mouth, or buccal mucosa, representing premalignant dysplastic changes.
    \item \textbf{Erythematous or wipeable plaques:} Creamy white patches that can be scraped off to reveal a raw, bleeding surface, characteristic of pseudomembranous oral candidiasis.
    \item \textbf{Salivary gland swelling:} Tender, unilateral swelling of the parotid or submandibular glands, particularly during meals, which may be accompanied by pus discharge from the duct orifice.
    \item \textbf{Glossodynia and altered taste:} A chronic burning sensation in the mouth (burning mouth syndrome) or dysgeusia (foul or metallic taste).
\end{itemize}

\section*{Differential Diagnosis}
It is critical to differentiate benign, self-limiting oral conditions from premalignant lesions, malignancies, and systemic diseases.
\begin{itemize}
    \item \textbf{Oral squamous cell carcinoma vs.\ Major aphthous ulcer:} Both present as large ulcers, but carcinoma is typically painless initially, indurated, has raised borders, and persists longer than 2--3 weeks, whereas aphthous ulcers are highly painful and resolve.
    \item \textbf{Oral lichen planus vs.\ Leukoplakia:} Lichen planus often presents with a characteristic bilateral, reticular (lace-like) pattern (Wickham's striae) on the buccal mucosa, whereas leukoplakia is typically a solitary, asymmetrical plaque.
    \item \textbf{Oral candidiasis vs.\ Oral hairy leukoplakia:} Hairy leukoplakia (associated with Epstein-Barr virus in immunocompromised patients) presents on the lateral borders of the tongue and cannot be scraped off, unlike candidiasis.
    \item \textbf{Pemphigus vulgaris vs.\ Mucous membrane pemphigoid:} Both are rare autoimmune blistering diseases, but pemphigus involves fragile intraepithelial blisters that rupture immediately, while pemphigoid features subepithelial blisters that are more stable.
\end{itemize}

\section*{When to Seek Medical Care}
Patients should seek medical or dental advice if they develop any new oral lesion, white or red patch, unexplained salivary swelling, or difficulty swallowing.

\textbf{Call 000 (or local emergency services) for:}
\begin{itemize}
    \item Rapidly progressive swelling of the neck, jaw, or tongue that compromises breathing or causes severe difficulty swallowing.
    \item Severe, acute facial swelling accompanied by high fever, confusion, or difficulty moving the eyes.
    \item Uncontrolled oral haemorrhage following trauma or surgery.
\end{itemize}

\section*{Diagnosis and Investigations}
Accurate diagnosis combines thorough clinical examination, tissue biopsy, and microbiological or imaging studies.
\begin{itemize}
    \item \textbf{Incisional or excisional biopsy:} The gold standard investigation for any chronic mucosal lesion persisting for more than 14 days, particularly leukoplakia, erythroplakia, or indurated ulcers, to rule out dysplasia or malignancy.
    \item \textbf{Exfoliative cytology and culture:} Swabbing lesions to identify fungal hyphae (Candida) or viral DNA (Herpes Simplex Virus via PCR).
    \item \textbf{Salivary gland imaging:} Ultrasound, sialography, or CT scanning to identify salivary duct stones (sialolithiasis) or evaluate glandular masses.
    \item \textbf{Autoimmune serology:} Testing for antinuclear antibodies (ANA), anti-SSA, and anti-SSB antibodies in suspected cases of Sjögren's syndrome or lupus.
    \item \textbf{Haematological investigations:} Full blood count, iron studies, folate, and vitamin B12 levels to investigate recurrent mouth ulcers or burning mouth symptoms.
\end{itemize}

\section*{Management}
Management is highly specific to the underlying disease, incorporating pharmacological, surgical, and supportive treatments.
\begin{itemize}
    \item \textbf{Antimicrobial therapy:} Topical nystatin or oral fluconazole for oral candidiasis; oral acyclovir for severe primary herpes simplex infections; and systemic antibiotics (such as cephalexin) for bacterial sialadenitis.
    \item \textbf{Anti-inflammatory and immunosuppressive therapies:} Topical corticosteroid gels (such as triamcinolone acetonide) for oral lichen planus and aphthous ulcers; systemic immunosuppression (prednisolone, azathioprine) for severe autoimmune conditions.
    \item \textbf{Surgical excision:} Wide local excision with neck dissection for oral squamous cell carcinoma, often combined with radiotherapy or chemotherapy.
    \item \textbf{Salivary gland interventions:} Sialolithotomy (surgical removal of salivary stones) or sialendoscopy to dilate ductal strictures.
    \item \textbf{Symptomatic relief:} Using topical anaesthetic mouthwashes (containing lidocaine or benzydamine) and artificial saliva sprays for xerostomia.
\end{itemize}

\section*{Complications}
Untreated or advanced mouth diseases can lead to severe localised tissue destruction, systemic spread of infection, and death.
\begin{itemize}
    \item \textbf{Metastasis of oral cancer:} Spread of oral squamous cell carcinoma to regional cervical lymph nodes and distant organs, significantly reducing survival rates.
    \item \textbf{Deep neck space infections:} Extension of salivary or dental infections into the retropharyngeal or submandibular spaces, leading to airway obstruction or mediastinitis.
    \item \textbf{Severe malnutrition and cachexia:} Chronic, severe oral pain (from ulcers, mucositis, or cancer) that prevents food intake, leading to rapid weight loss and nutritional depletion.
    \item \textbf{Systemic sepsis:} Spread of oral infections in severely immunocompromised individuals, potentially causing life-threatening systemic blood infections.
\end{itemize}

\section*{Prognosis and Outcomes}
The prognosis depends on the specific condition, but is heavily influenced by the timeliness of diagnosis. Benign conditions like aphthous ulcers and oral candidiasis have an excellent prognosis, resolving within days to weeks with minimal sequelae. Autoimmune conditions like oral lichen planus are chronic and require long-term management, but are generally well-controlled with topical therapies. For oral squamous cell carcinoma, the prognosis is highly stage-dependent: the 5-year survival rate is approximately 80--90\% for early-stage localised disease, but drops to less than 30--40\% if regional lymph node metastasis has occurred, highlighting the critical importance of early biopsy and detection.

\section*{Prevention and Self-Care}
Preventing mouth diseases involves minimising exposure to carcinogens, practising good hygiene, and performing regular self-examinations.
\begin{itemize}
    \item \textbf{Perform monthly oral self-exams:} Inspect the lips, tongue, cheeks, and floor of the mouth in a well-lit mirror, looking for any new red or white patches, lumps, or non-healing ulcers.
    \item \textbf{Quit tobacco and limit alcohol:} Stop smoking or chewing tobacco, and moderate alcohol consumption, as their synergistic effect is the leading cause of oral cancer.
    \item \textbf{Practise strict oral hygiene:} Brush and floss daily, and clean dentures thoroughly every night to prevent denture stomatitis (candidiasis).
    \item \textbf{Stay hydrated:} Drink plenty of water throughout the day, and chew sugar-free gum to stimulate salivary flow and protect the oral mucosa.
    \item \textbf{Protect lips from UV radiation:} Apply a lip balm with SPF 30 or higher when outdoors to prevent actinic cheilitis (premalignant lip changes caused by sun exposure).
\end{itemize}

\section*{Sources and Further Reading}
The following reputable organisations provide further information and clinical guidelines on mouth diseases and conditions:
\begin{itemize}
    \item Australasian College of Dermatologists. \textit{Oral mucosal diseases.} https://www.dermcoll.edu.au/
    \item Better Health Channel. \textit{Mouth conditions.} https://www.betterhealth.vic.gov.au/
    \item Cancer Council Australia. \textit{Oral cancer.} https://www.cancer.org.au/
    \item DermNet. \textit{Oral mucosal lesions.} https://dermnetnz.org/
    \item Mayo Clinic. \textit{Mouth diseases and conditions.} https://www.mayoclinic.org/
\end{itemize}
